\documentclass[11pt]{article}

\usepackage[T1]{fontenc}
\usepackage[utf8]{inputenc}
\usepackage{lmodern}
\usepackage{amsmath,amssymb,amsfonts,amsthm}
\usepackage[margin=1in]{geometry}
\usepackage{microtype}
\usepackage{xcolor}
\usepackage[colorlinks=true,linkcolor=teal,urlcolor=teal]{hyperref}

\newcommand{\Hc}{\mathsf H}
\newcommand{\F}{\mathcal F_{\mathrm{orb}}}
\newcommand{\T}{\mathcal T_1}
\newcommand{\Tr}{\operatorname{Tr}}
\newcommand{\dd}{\,\mathrm d}
\newcommand{\B}{\mathrm B}
\newcommand{\Bcl}{B_{\mathrm{cl}}}
\newcommand{\Borb}{B_{\mathrm{orb}}}
\newcommand{\Buniv}{B_{\mathrm{univ}}}

\newtheorem*{propC}{Proposition C.6 (rephrased)}

\title{A Short Guide to Proposition C.6\\
\large The hybrid Gaussian flat-prior upper bound}
\author{}
\date{}

\begin{document}
\maketitle

\begin{abstract}
Proposition C.6 is the converse for the hybrid Gaussian experiment containing
both the classical spectral variable and the quantum orbital modes.  It
shows that even a channel coupling these two components cannot beat the
product value
$\Buniv(\gamma,p)=\Bcl(\gamma,d)\Borb(\gamma,p)$.  The proof uses a
classical Gaussian witness and a compact quantum witness, extracts a
covariant quantum limit and a translation-covariant scalar limit, and bounds
them separately.  This note organizes the argument into its two parallel
parts and explains how they recombine into the sharp fidelity bound.
\end{abstract}

\section{The hybrid Gaussian experiment}

The classical space is
\[
  \Hc=\left\{h\in\mathbb R^d:\sum_i h_i=0\right\},
  \qquad \dim\Hc=d-1,
\]
with positive covariance $\Sigma_p$.  Write
\[
  f_1=\varphi_{0,\Sigma_p},
  \qquad
  f_\gamma=\varphi_{0,\gamma\Sigma_p}
\]
for the centered Gaussian densities.  The orbital component lives on the
multimode Fock space $\F$ and has centered thermal state $\tau$.

At parameter $\xi=(h,z)\in\Hc\times\mathbb C^{r_d}$, the input and ideal
target are
\begin{align*}
  \Theta_{h,z}^{\mathrm{in}}
  &=\mathsf N_{h,\Sigma_p}\otimes\Phi_z,\\
  \Theta_{h,z}^{\mathrm{tar}}
  &=\mathsf N_{\sqrt\gamma h,\Sigma_p}\otimes\Psi_z.
\end{align*}
The target amplifies both means by $\sqrt\gamma$ without increasing either
the classical covariance or the centered thermal noise.  A physical
procedure cannot do this perfectly.

The optimal product procedure dilates the classical observation by
$\sqrt\gamma$ and applies the product quantum-limited amplifier.  Its two
losses are
\begin{align}
  \Bcl(\gamma,d)
  &=\left(\frac{2\sqrt\gamma}{1+\gamma}\right)^{(d-1)/2},
  \label{eq:Bcl}\\
  \Borb(\gamma,p)
  &=\B(\tau^{(\gamma)},\tau),
  \label{eq:Borb}
\end{align}
and root fidelity multiplies across the classical and quantum factors.
Proposition C.6 proves that a joint hybrid channel cannot improve on their
product.

\section{Statement of Proposition C.6}

Let $\Omega_L$ be the centered cube in
$\Hc\times\mathbb C^{r_d}$ used in the manuscript.

\begin{propC}
For the hybrid Gaussian input and target families,
\[
  \limsup_{L\to\infty}\sup_\Lambda
  \frac1{|\Omega_L|}\int_{\Omega_L}
  \B\!\left(\Lambda(\Theta_{h,z}^{\mathrm{in}}),
  \Theta_{h,z}^{\mathrm{tar}}\right)\dd h\dd z
  \leq\Buniv(\gamma,p),
\]
where
$\Buniv(\gamma,p)=\Bcl(\gamma,d)\Borb(\gamma,p)$ and the supremum is over
all hybrid channels, including channels that couple the classical and
quantum components.
\end{propC}

\section{The two witnesses}

\subsection{Classical witness}

Define the square-root likelihood ratio
\begin{equation}\label{eq:chi}
  \chi(y)=\sqrt{\frac{f_1(y)}{f_\gamma(y)}}
  =\gamma^{(d-1)/4}
  \exp\!\left[-\frac{\gamma-1}{4\gamma}
  \langle y,\Sigma_p^{-1}y\rangle\right].
\end{equation}
It is bounded, strictly positive, and belongs to $C_0(\Hc)$.  Moreover,
\begin{equation}\label{eq:classical-id}
  \int_{\Hc}\frac{f_1(y)}{\chi(y)}\dd y
  =\int_{\Hc}\sqrt{f_1(y)f_\gamma(y)}\dd y
  =\Bcl(\gamma,d).
\end{equation}
The fact that $\chi$ vanishes at infinity is as important as the Gaussian
identity: it makes integration against $\chi$ continuous for weak-* limits
of finite measures.

\subsection{Quantum witness}

In the joint number basis of $\tau$ and the quantum-limited output
$\tau^{(\gamma)}$, define
\begin{equation}\label{eq:W}
  W_\gamma
  =\sum_{\mathbf k}
  \sqrt{\frac{t_{\mathbf k}}{t_{\mathbf k}^{(\gamma)}}}
  |\mathbf k\rangle\!\langle\mathbf k|.
\end{equation}
It is positive, injective, bounded, compact, and a product of decreasing
photon-number weights.  Its matching identities are
\begin{equation}\label{eq:quantum-id}
  \Tr(\tau W_\gamma^{-1})
  =\Tr(\tau^{(\gamma)}W_\gamma)
  =\Borb(\gamma,p).
\end{equation}

Lemma C.1 combines these witnesses.  If $R(y)$ is a positive
classical--quantum density, then
\begin{equation}\label{eq:weighted-hybrid}
  \B(R,f_1\tau)^2
  \leq
  \left(\int_{\Hc}\chi(y)\Tr[R(y)W_\gamma]\dd y\right)
  \Bcl(\gamma,d)\Borb(\gamma,p).
\end{equation}
It remains to show that the first, channel-dependent factor is at most the
same product $\Bcl\Borb$.

\section{Proof roadmap}

The proof performs the following reductions:
\[
\begin{array}{c}
\text{average fidelity over a large parameter cube}\\
\Downarrow\\
\text{fidelity at the origin for a F\o lner-averaged channel}\\
\Downarrow\\
\text{one weighted moment }M_j\\
\Downarrow\\
\begin{array}{c@{\qquad}c}
\text{quantum compact limit }\Gamma^\chi
&\text{classical scalar limit }\mathcal R\\
\text{Lemma C.2}&\text{Lemma C.5}
\end{array}\\
\Downarrow\\
M_j\leq\Bcl(\gamma,d)\Borb(\gamma,p)\\
\Downarrow\\
\B\leq\Bcl(\gamma,d)\Borb(\gamma,p).
\end{array}
\]
The two limits come from the same averaged channels but retain different
information.  The compact quantum limit controls the $W_\gamma$ moment; the
scalar limit controls how much classically weighted trace that quantum map
can have.

\section{Proof walkthrough}

\subsection*{Step 1: F\o lner averaging}

For $\xi=(h,z)$, let
\begin{align*}
  (\mathcal U_\xi^{\mathrm{in}}R)(x)
  &=D(z)R(x-h)D(z)^*,\\
  (\mathcal U_\xi^{\mathrm{out}}R)(y)
  &=D(\sqrt\gamma z)R(y-\sqrt\gamma h)D(\sqrt\gamma z)^*.
\end{align*}
Then
\[
  \Theta_\xi^{\mathrm{in}}
  =\mathcal U_\xi^{\mathrm{in}}(f_1\tau),
  \qquad
  \Theta_\xi^{\mathrm{tar}}
  =\mathcal U_\xi^{\mathrm{out}}(f_1\tau).
\]
For a hybrid channel $\Lambda$, define
\begin{equation}\label{eq:average}
  \overline\Lambda_L
  =\frac1{|\Omega_L|}\int_{\Omega_L}
  \mathcal U_{-\xi}^{\mathrm{out}}\circ\Lambda\circ
  \mathcal U_\xi^{\mathrm{in}}\dd\xi.
\end{equation}
Joint concavity and invariance of root fidelity give
\begin{equation}\label{eq:concavity}
  \frac1{|\Omega_L|}\int_{\Omega_L}
  \B(\Lambda(\Theta_\xi^{\mathrm{in}}),
    \Theta_\xi^{\mathrm{tar}})\dd\xi
  \leq\B(\overline\Lambda_L(f_1\tau),f_1\tau).
\end{equation}

Choose $L_j\to\infty$ and channels approaching the limsup, and denote their
averages by $\Lambda_j$.  Translating the averaging cube shows that for each
fixed $\xi$,
\begin{equation}\label{eq:approx-covariance}
  \|\Lambda_j(\mathcal U_\xi^{\mathrm{in}}R)
  -\mathcal U_\xi^{\mathrm{out}}\Lambda_j(R)\|_1
  \leq
  \frac{|(\Omega_{L_j}+\xi)\mathbin\triangle\Omega_{L_j}|}
       {|\Omega_{L_j}|}\,\|R\|_1
  \longrightarrow0.
\end{equation}

\subsection*{Step 2: reduce fidelity to one moment}

Put
\[
  R_j=\Lambda_j(f_1\tau),
  \qquad
  M_j=\int_{\Hc}\chi(y)\Tr[R_j(y)W_\gamma]\dd y.
\]
The weighted hybrid inequality \eqref{eq:weighted-hybrid} gives
\begin{equation}\label{eq:B-to-M}
  \B(R_j,f_1\tau)^2
  \leq M_j\,\Bcl(\gamma,d)\Borb(\gamma,p).
\end{equation}
The rest of the proof is devoted to the upper bound
\begin{equation}\label{eq:M-goal}
  \limsup_j M_j
  \leq\Bcl(\gamma,d)\Borb(\gamma,p).
\end{equation}

\subsection*{Step 3: extract the weighted quantum map}

For a trace-class operator $X$ on $\F$, define
\begin{equation}\label{eq:Gamma-j}
  \Gamma_j^\chi(X)
  =\int_{\Hc}\chi(y)[\Lambda_j(f_1X)](y)\dd y.
\end{equation}
This map feeds $\Lambda_j$ a fixed classical Gaussian with quantum input
$X$, then averages the quantum output with the classical weight $\chi$.
It is completely positive and satisfies
\[
  \Tr\Gamma_j^\chi(X)
  \leq\|\chi\|_\infty\Tr X
  \qquad(X\geq0).
\]
Notice that $\|\chi\|_\infty=\gamma^{(d-1)/4}$ may exceed one.  Hence this
initial estimate does \emph{not} make $\Gamma_j^\chi$ trace nonincreasing.

Taking $h=0$ in \eqref{eq:approx-covariance} shows that these maps become
displacement covariant in the quantum variable.  Lemma C.2 therefore gives,
after taking a subsequence, a completely positive exactly covariant map
$\Gamma^\chi$ with
\begin{equation}\label{eq:quantum-limit}
  \Tr[\Gamma_j^\chi(X)K]
  \longrightarrow\Tr[\Gamma^\chi(X)K]
  \qquad(K\text{ compact}).
\end{equation}
Because $W_\gamma$ is compact,
\begin{equation}\label{eq:M-limit}
  M_j=\Tr[\Gamma_j^\chi(\tau)W_\gamma]
  \longrightarrow\Tr[\Gamma^\chi(\tau)W_\gamma].
\end{equation}

To apply the least-noise moment to $\Gamma^\chi$, the proof still has to
improve its trace bound from $\|\chi\|_\infty$ to one.  This is the purpose
of the classical limit.

\subsection*{Step 4: extract the scalar trace kernel}

For $u\in L^1(\Hc)$, let $\mathcal R_j(u)$ be the classical trace measure of
$\Lambda_j(u\tau)$:
\begin{equation}\label{eq:Rj}
  \mathcal R_j(u)(E)
  =\int_E\Tr[\Lambda_j(u\tau)(y)]\dd y.
\end{equation}
Since $\Lambda_j$ is a channel, $\mathcal R_j$ is a positive contraction
from $L^1(\Hc)$ to the finite Radon measures.

A diagonal extraction on countable dense subsets of $L^1(\Hc)$ and
$C_0(\Hc)$ gives a positive contraction $\mathcal R$ such that
\[
  \mathcal R_j(u)\stackrel{w^*}{\longrightarrow}\mathcal R(u)
  \qquad(u\in L^1(\Hc)).
\]
Taking $z=0$ in \eqref{eq:approx-covariance} and passing to this limit yields
\begin{equation}\label{eq:R-covariance}
  \mathcal R(u(\,\cdot-h))(E)
  =\mathcal R(u)(E-\sqrt\gamma h).
\end{equation}
Lemma C.5 identifies this covariant scalar kernel: there is a finite positive
measure $\mu$, with $\mu(\Hc)\leq1$, such that
\begin{equation}\label{eq:convolution}
  \mathcal R(u)(E)
  =\int_{\Hc}u(x)\mu(E-\sqrt\gamma x)\dd x.
\end{equation}

\subsection*{Step 5: bound the trace of the weighted quantum limit}

Let $P_N\uparrow I$ be finite-rank projections on $\F$.  Positivity and
\eqref{eq:quantum-limit} give
\begin{align}
  \Tr\Gamma^\chi(\tau)
  &=\sup_N\Tr[\Gamma^\chi(\tau)P_N]\notag\\
  &\leq\liminf_j\Tr\Gamma_j^\chi(\tau).
  \label{eq:trace-liminf}
\end{align}
By \eqref{eq:Gamma-j} and \eqref{eq:Rj},
\[
  \Tr\Gamma_j^\chi(\tau)
  =\int_{\Hc}\chi(y)\mathcal R_j(f_1)(\dd y).
\]
Because $\chi\in C_0(\Hc)$, weak-* convergence of the scalar measures passes
this integral to the limit.  Hence
\begin{equation}\label{eq:trace-to-classical}
  \Tr\Gamma^\chi(\tau)
  \leq\int_{\Hc}\chi(y)\mathcal R(f_1)(\dd y).
\end{equation}

Use the convolution representation \eqref{eq:convolution}.  If
$X\sim\mathsf N_{0,\Sigma_p}$, then $\sqrt\gamma X$ has density $f_\gamma$,
so
\begin{align}
  \int_{\Hc}\chi(y)\mathcal R(f_1)(\dd y)
  &=\int_{\Hc}c_\gamma(r)\mu(\dd r),
  \label{eq:cgamma-average}\\
  c_\gamma(r)
  &=\int_{\Hc}\chi(y)f_\gamma(y-r)\dd y.
  \label{eq:cgamma}
\end{align}
Completing the Gaussian square gives
\begin{equation}\label{eq:cgamma-value}
  c_\gamma(r)
  =\Bcl(\gamma,d)
  \exp\!\left[-\frac{\gamma-1}{2\gamma(\gamma+1)}
    \langle r,\Sigma_p^{-1}r\rangle\right]
  \leq\Bcl(\gamma,d).
\end{equation}
The maximum occurs at zero added classical noise.  Since
$\mu(\Hc)\leq1$, equations
\eqref{eq:trace-to-classical}--\eqref{eq:cgamma-value} yield
\begin{equation}\label{eq:trace-bound}
  \Tr\Gamma^\chi(\tau)\leq\Bcl(\gamma,d).
\end{equation}

This estimate is first proved only at $\tau$.  Exact displacement covariance
makes it global.  Indeed, the positive functional
$X\mapsto\Tr\Gamma^\chi(X)$ is represented by a bounded positive operator
$A$.  Covariance implies that $A$ commutes with every multimode Weyl
displacement, so irreducibility gives $A=cI$.  Since $\Tr\tau=1$,
\[
  c=\Tr\Gamma^\chi(\tau)leq\Bcl(\gamma,d)\leq1.
\]
Therefore $\Gamma^\chi$ is trace nonincreasing and Proposition C.3 applies.

\subsection*{Step 6: apply the quantum least-noise bound}

Using \eqref{eq:M-limit}, Proposition C.3, and
\eqref{eq:quantum-id},
\begin{align*}
  \lim_j M_j
  &=\Tr[\Gamma^\chi(\tau)W_\gamma]\\
  &\leq\Tr\Gamma^\chi(\tau)
    \Tr(\tau^{(\gamma)}W_\gamma)\\
  &\leq\Bcl(\gamma,d)\Borb(\gamma,p).
\end{align*}
This proves \eqref{eq:M-goal}.  Substitution into \eqref{eq:B-to-M} gives
\[
  \limsup_j\B(R_j,f_1\tau)^2
  \leq
  \bigl(\Bcl(\gamma,d)\Borb(\gamma,p)\bigr)^2.
\]
Taking square roots and using the F\o lner concavity bound
\eqref{eq:concavity} proves Proposition C.6.

\section{Why the classical and quantum factors multiply}

The proof does not assume that the competing channel factors.  Instead,
factorization emerges at the level of witnesses:
\begin{itemize}
  \item $\chi$ is the square-root likelihood ratio for the two classical
  Gaussian covariances;
  \item $W_\gamma$ is the square-root likelihood-ratio observable for the two
  thermal laws;
  \item Lemma C.1 separates the channel-dependent moment from the two fixed
  inverse-witness integrals;
  \item Lemma C.5 bounds the trace available to the quantum moment by the
  classical affinity $\Bcl$;
  \item Proposition C.3 bounds the normalized quantum moment by the orbital
  affinity $\Borb$.
\end{itemize}
Thus any classical--quantum coupling is already included in the single
moment $M_j$, yet the covariance constraints force that moment to pay both
losses.

\section{Common points of confusion}

\begin{enumerate}
  \item The compact-limit lemma alone gives
  $\Tr\Gamma^\chi(X)\leq\|\chi\|_\infty\Tr X$, which is not enough for
  Proposition C.3.  The scalar convolution calculation is what improves the
  coefficient to $\Bcl\leq1$.
  \item Two different compactness arguments are used.  Lemma C.2 treats
  quantum maps through compact operators, while the diagonal extraction of
  $\mathcal R_j$ treats classical finite measures through $C_0(\Hc)$.
  \item Compactness of $W_\gamma$ passes the quantum moment to the limit;
  membership $\chi\in C_0(\Hc)$ passes the classical weighted trace to the
  limit.  These are parallel requirements.
  \item The noise measure $\mu$ need not be normalized or absolutely
  continuous.  The estimate uses only positivity and
  $\mu(\Hc)\leq1$.
\end{enumerate}

\section{Role in the paper}

Proposition C.6 is the upper-bound half of the hybrid Gaussian flat-prior
theorem.  The product of classical dilation and quantum-limited amplification
attains the same value, so
\[
  \Buniv(\gamma,p)
  =\Bcl(\gamma,d)\Borb(\gamma,p)
\]
is sharp.  LAN fidelity transfer then carries this Gaussian converse back to
the unknown-spectrum cloning problem and produces the classical penalty
absent from known-spectrum cloning.

\end{document}
